\chapter{Tag 5 — Eine neue Mathematik: endliche Körper}

\begin{quote}
Gestern hast du mit CRC etwas Erstaunliches getan: du hast \textbf{mit
Polynomen gerechnet, modulo 2}. Heute gehen wir einen Schritt weiter
und entdecken, dass dahinter eine ganze mathematische Welt steckt:
die der \textbf{endlichen Körper}. Das klingt sperrig, ist aber im
Grunde nur „Modulo-Rechnung, hochgezogen auf Polynome“. Und genau
diese Welt brauchen wir, um morgen Reed-Solomon zu verstehen — das
Verfahren, das im Datamatrix-Code steckt.
\end{quote}

\section*{Lernziele für heute}

Am Ende von Tag 5 kannst du \dots

\begin{itemize}
\item erklären, was ein \textbf{Körper} in der Mathematik ist und
  warum „Körper“ hier nicht „Volumen“ heißt
\item die Rechenregeln in \textbf{GF(2)} (dem kleinsten Körper)
  anwenden
\item zeigen, warum \textbf{Modulo 256} für Bytes nicht funktioniert
\item in \textbf{GF($2^3$)} addieren und multiplizieren — mit
  Bleistift und Zettel
\item für jedes Element in GF($2^3$) das \textbf{multiplikative
  Inverse} finden
\item in Python eine kleine Klasse für GF($2^3$) bauen und damit
  experimentieren
\end{itemize}

Material: Bleistift, kariertes Papier, viel Geduld für die
Multiplikationstafel. Optional: zwei verschiedene Farben (für
Datenbits und Polynom-Koeffizienten).

% =============================================================
\section{Was ist eigentlich ein „Körper“? \normalfont(\textit{≈ 25 min})}
% =============================================================

Bisher haben wir mit ganz unterschiedlichen Zahlensystemen gerechnet:
ganze Zahlen, modulo 10 (EAN-13), modulo 11 (ISBN), modulo 2
(Parität), Polynome modulo 2 (CRC). Lass uns mal sortieren, was die
alle gemeinsam haben.

In jedem dieser Systeme kannst du:

\begin{itemize}
\item \textbf{Addieren} (immer)
\item \textbf{Subtrahieren} (immer; bei Modulo wird's „addieren mit
  Gegenzahl“)
\item \textbf{Multiplizieren} (immer)
\item \textbf{Dividieren} — das ist die spannende Frage!
\end{itemize}

\subsection*{Wann kann man dividieren?}

Dividieren heißt: zu jeder Zahl $\ne 0$ gibt es eine \textbf{Gegenzahl}
(das \emph{multiplikative Inverse}), sodass $a \cdot a^{-1} = 1$. Bei
den reellen Zahlen ist das einfach: $5^{-1} = 0{,}2$. Bei
modulo-Systemen wird's interessant.

\begin{bleistiftuebung}{1}{Wann gibt es Inverse?}
\begin{enumerate}[label=(\alph*), leftmargin=*]
\item Versuche, in \textbf{modulo 10} für jede Zahl 1 bis 9 ein
  Inverses zu finden, also eine Zahl $b$, sodass $a \cdot b \equiv 1
  \pmod{10}$. Welche Zahlen haben Inverse, welche nicht?

  \begin{quote}
  Tipp: probiere systematisch. Für $a = 3$ z.\,B.\ die Reihe
  $3 \cdot 1 = 3,\ 3 \cdot 2 = 6,\ 3 \cdot 3 = 9,\ 3 \cdot 4 = 12
  \equiv 2, \dots$ und schaue, wann die 1 auftaucht.
  \end{quote}

\item Mache dasselbe für \textbf{modulo 11}. Welche Zahlen haben
  Inverse?

\item Was fällt dir auf? Welche Eigenschaft hat 11, die 10 nicht hat?

\item \textbf{Faustregel formulieren:} bei Modulo $n$ hat eine Zahl
  $a$ genau dann ein Inverses, wenn \dots?
\end{enumerate}
\end{bleistiftuebung}

\subsection*{Was ist nun ein Körper?}

Ein \textbf{Körper} (engl.\ \emph{field}) ist eine Menge von Zahlen,
in der man \textbf{uneingeschränkt rechnen} kann: $+, -, \cdot, \div$
(außer Division durch 0). Konkret heißt das:

\begin{itemize}
\item Es gibt eine 0 (neutrales Element der Addition).
\item Es gibt eine 1 (neutrales Element der Multiplikation).
\item Jede Zahl hat eine Gegenzahl (additives Inverses).
\item Jede Zahl $\ne 0$ hat ein multiplikatives Inverses.
\item Plus die üblichen Rechenregeln (Assoziativität, Kommutativität,
  Distributivität).
\end{itemize}

Beispiele für Körper:
\begin{itemize}
\item $\mathbb{R}$ (reelle Zahlen)
\item $\mathbb{Q}$ (rationale Zahlen)
\item $\mathbb{C}$ (komplexe Zahlen)
\item $\mathbb{Z}/p\mathbb{Z}$ — die ganzen Zahlen modulo $p$ mit $p$
  prim (das ist neu!)
\end{itemize}

Beispiele für \textbf{keine} Körper:
\begin{itemize}
\item $\mathbb{Z}$ (ganze Zahlen) — kein Inverses für 2 (es gibt kein
  „$1/2$“ in $\mathbb{Z}$)
\item $\mathbb{Z}/10\mathbb{Z}$ — modulo 10, weil 2, 4, 5, 6, 8 keine
  Inverse haben
\end{itemize}

\begin{quote}
Merke: \textbf{Körper modulo $n$ gibt es nur, wenn $n$ eine Primzahl
ist.} Diese Körper schreibt man oft auch GF($p$) — „Galois-Feld“ der
Größe $p$, benannt nach dem französischen Mathematiker Évariste Galois
(1811–1832, lebte nur 20 Jahre, hat aber eine ganze Disziplin der
Mathematik begründet).
\end{quote}

\begin{bleistiftuebung}{2}{Der kleinste Körper: GF(2)}
GF(2) ist der allerkleinste Körper. Er hat genau zwei Elemente: 0 und
1. Du kennst ihn schon — nur unter einem anderen Namen.

\begin{enumerate}[label=(\alph*), leftmargin=*]
\item Erstelle die \textbf{Additionstafel} für GF(2):

  \begin{center}
  \begin{tabular}{c|cc}
    $+$ & 0 & 1 \\
    \hline
    0 & & \\
    1 & & \\
  \end{tabular}
  \end{center}

\item Erstelle die \textbf{Multiplikationstafel} für GF(2):

  \begin{center}
  \begin{tabular}{c|cc}
    $\cdot$ & 0 & 1 \\
    \hline
    0 & & \\
    1 & & \\
  \end{tabular}
  \end{center}

\item Vergleiche: welche aus Python bekannten Operationen ergeben
  dieselben Tabellen?

  \begin{quote}
  Tipp: erinnere dich an Tag 1 (Parität) und Tag 4 (CRC).
  \end{quote}

\item Was ist $1 - 1$ in GF(2)? Was ist also die Gegenzahl von 1?
  (Gemeinheit: in GF(2) ist das nicht „$-1$“.)
\end{enumerate}
\end{bleistiftuebung}

% =============================================================
\section{Warum Modulo 256 nicht der Weg ist \normalfont(\textit{≈ 15 min})}
% =============================================================

Wir wollen mit Bytes rechnen, also mit Zahlen 0 bis 255. Naheliegend
wäre: einfach Modulo 256. Aber 256 ist keine Primzahl ($256 = 2^8$),
also haben wir das Problem aus Bleistiftübung 1 zurück.

\begin{bleistiftuebung}{3}{Modulo 256 ist kaputt}
\begin{enumerate}[label=(\alph*), leftmargin=*]
\item Was ist $2 \cdot 128 \bmod 256$? Folgerung: hat 2 ein Inverses
  in modulo 256?

\item Welche Zahlen aus $1 \dots 255$ haben kein Inverses modulo 256?

  \begin{quote}
  Tipp: eine Zahl $a$ hat genau dann ein Inverses modulo $n$, wenn
  $a$ und $n$ keinen gemeinsamen Teiler außer 1 haben (man sagt:
  \emph{teilerfremd} oder \emph{koprim}).
  \end{quote}

\item Wie viele Zahlen aus $1 \dots 255$ sind koprim zu 256? (Du
  musst nicht alle aufzählen, aber überlege, welche es nicht sind.)
\end{enumerate}
\end{bleistiftuebung}

\subsection*{Zwei mögliche Auswege}

Wir hätten zwei Optionen:

\textbf{Option 1:} statt 256 eine Primzahl in der Nähe nehmen,
z.\,B.\ 257 (das ist tatsächlich prim!). Damit hätten wir GF(257) —
einen schönen Körper. Aber: Bytes sind 256 Werte, nicht 257. Wir
würden ständig den 257-ten Wert „übrig“ haben. Unschön.

\textbf{Option 2:} eine ganz neue Konstruktion: statt mit Zahlen
modulo $2^8$ rechnen wir mit \textbf{Polynomen} vom Grad $< 8$ über
GF(2), und nehmen modulo eines bestimmten Polynoms. Klingt fies, ist
aber elegant — und genau das, was wir heute aufbauen.

Wir starten klein: erst mit GF($2^3$) (8 Elemente), dann später mit
GF($2^8$) (256 Elemente).

% =============================================================
\section{GF($2^3$) konkret aufbauen \normalfont(\textit{≈ 60 min})}
% =============================================================

Das ist der schwierigste Teil heute. Aber wenn du ihn verstanden
hast, hast du das Tor zu Reed-Solomon weit aufgestoßen.

\subsection*{Die Grundidee}

Wir betrachten \textbf{Polynome vom Grad kleiner als 3, deren
Koeffizienten 0 oder 1 sind}. Das sind genau $2^3 = 8$ Stück:

\begin{verbatim}
0
1
x
x+1
x²
x²+1
x²+x
x²+x+1
\end{verbatim}

Diese 8 Polynome sind die Elemente von GF($2^3$). Jedes lässt sich
als 3-Bit-Zahl schreiben:

\begin{verbatim}
0       =  000
1       =  001
x       =  010
x+1     =  011
x²      =  100
x²+1    =  101
x²+x    =  110
x²+x+1  =  111
\end{verbatim}

Praktischer Trick: wir können also einfach mit 3-Bit-Zahlen rechnen,
müssen aber die Rechenregeln der Polynome beachten.

\subsection*{Addition in GF($2^3$)}

Polynome addieren wir koeffizientenweise. Da die Koeffizienten in
GF(2) liegen (also 0 oder 1), ist die Addition modulo 2 — genau
\textbf{XOR}.

Beispiel: $(x^2+x+1) + (x^2+1)$. Koeffizienten:
\begin{verbatim}
   1 1 1
⊕  1 0 1
─────────
   0 1 0   →  x
\end{verbatim}

Also: $(x^2+x+1) + (x^2+1) = x$. Schreibst du das als 3-Bit-Zahlen:
$\texttt{111} \oplus \texttt{101} = \texttt{010}$.

\begin{bleistiftuebung}{4}{Addition in GF($2^3$)}
Berechne (am einfachsten als XOR der 3-Bit-Darstellungen):

\begin{center}
\begin{tabular}{lll}
  \toprule
  Aufgabe & Ergebnis (3-Bit) & Als Polynom \\
  \midrule
  $(x+1) + (x^2+x)$            & & \\
  $(x^2+1) + (x^2+x+1)$        & & \\
  $x^2 + x^2$                  & & \\
  $(x^2+x+1) + (x^2+x+1)$      & & \\
  \bottomrule
\end{tabular}
\end{center}

\begin{quote}
Beobachtung: was ist $a + a$ für jedes Element $a$?
\end{quote}
\end{bleistiftuebung}

\subsection*{Multiplikation — jetzt wird's interessant}

Polynome multiplizieren wir wie in der Schule, aber \textbf{die
Koeffizienten in GF(2)}. Beispiel:

\begin{verbatim}
(x + 1) · (x + 1)
= x·x + x·1 + 1·x + 1·1
= x² + x + x + 1
= x² + 0 + 1            (denn x + x = 0 in GF(2))
= x² + 1
\end{verbatim}

Aber: was, wenn das Ergebnis Grad $\ge 3$ hat? Dann passt es nicht
mehr in unsere 3-Bit-Welt.

\begin{verbatim}
(x²) · (x) = x³        ← Grad 3, zu hoch!
\end{verbatim}

\subsection*{Modulo-Polynom: das irreduzible Polynom}

Hier kommt der Schlüssel-Trick. So wie bei modulo $p$ alles, was
$\ge p$ wird, modulo $p$ reduziert wird, nehmen wir hier
\textbf{alles modulo eines bestimmten Polynoms} $p(x)$ mit Grad 3.

Welches Polynom nehmen wir? Es muss \textbf{irreduzibel} sein — das
heißt, es darf nicht in Faktoren niedrigeren Grades zerlegbar sein
(analog zu „Primzahl“ bei Zahlen).

Ein irreduzibles Polynom vom Grad 3 über GF(2) ist:
\[
  p(x) = x^3 + x + 1
\]
(Es gibt noch genau ein anderes, $x^3 + x^2 + 1$, aber das nehmen
wir nicht.)

\subsection*{Modulo $p(x)$ rechnen}

Wenn ein Multiplikationsergebnis ein $x^3$ enthält, ersetzen wir es:
aus $p(x) = x^3 + x + 1 = 0$ (in unserer neuen Welt!) folgt
$x^3 = x + 1$. Diese \textbf{Reduktionsregel} brauchen wir ständig:
\begin{align*}
  x^3  &=  x + 1 \\
  x^4  &=  x \cdot x^3 = x \cdot (x + 1) = x^2 + x \\
  x^5  &=  x \cdot x^4 = x \cdot (x^2 + x) = x^3 + x^2 = (x+1) + x^2 = x^2 + x + 1
\end{align*}

\begin{bleistiftuebung}{5}{Multiplikation in GF($2^3$)}
Berechne, mit $p(x) = x^3 + x + 1$:

\begin{enumerate}[label=(\alph*), leftmargin=*]
\item $x \cdot x^2$ (kleine Aufwärmübung)
\item $x^2 \cdot x^2$

  \begin{quote}
  Tipp: $x^2 \cdot x^2 = x^4$, und $x^4 = x^2 + x$ (siehe oben).
  \end{quote}

\item $(x+1) \cdot (x^2+x)$

\item $(x^2+x+1) \cdot (x^2+1)$

  \begin{quote}
  Tipp: erst auspolieren (Distributivgesetz), dann reduzieren.
  \end{quote}
\end{enumerate}
\end{bleistiftuebung}

\subsection*{Die vollständige Multiplikationstafel}

Wenn du Lust hast, kannst du dir die komplette $8 \times 8$-Multi\-pli\-ka\-tions\-tafel
von GF($2^3$) selbst ausrechnen. Das ist mühsam, aber sehr lehrreich
— und du machst es nur einmal.

\begin{bleistiftuebung}{6}{(Bonus) Die Tafel}
Fülle die Tafel aus. Element 0 lasse ich dir geschenkt, das ist die
ganze erste Zeile/Spalte.

\begin{center}\small
\begin{tabular}{l|lllllll}
  \toprule
  $\cdot$  & $1$ & $x$ & $x{+}1$ & $x^2$ & $x^2{+}1$ & $x^2{+}x$ & $x^2{+}x{+}1$ \\
  \midrule
  $1$            & $1$ & $x$ & $x{+}1$ & $x^2$ & $x^2{+}1$ & $x^2{+}x$ & $x^2{+}x{+}1$ \\
  $x$            & $x$ &     &         &       &           &           &               \\
  $x{+}1$        &     &     &         &       &           &           &               \\
  $x^2$          &     &     &         &       &           &           &               \\
  $x^2{+}1$      &     &     &         &       &           &           &               \\
  $x^2{+}x$      &     &     &         &       &           &           &               \\
  $x^2{+}x{+}1$  &     &     &         &       &           &           &               \\
  \bottomrule
\end{tabular}
\end{center}

\begin{quote}
Wenn du das durchhältst, hast du mehr über endliche Körper gelernt
als in den meisten Mathematik-Erstsemestern.
\end{quote}
\end{bleistiftuebung}

\subsection*{Inverse finden}

Sobald du die Multiplikationstafel hast, sind die Inverse einfach
abzulesen: das Inverse von $a$ ist das Element $b$, sodass
$a \cdot b = 1$.

\begin{bleistiftuebung}{7}{Inverse aus der Tafel}
Lies aus deiner Tafel ab:

\begin{center}
\begin{tabular}{ll}
  \toprule
  Element & Inverses \\
  \midrule
  $1$            & \\
  $x$            & \\
  $x+1$          & \\
  $x^2$          & \\
  $x^2+1$        & \\
  $x^2+x$        & \\
  $x^2+x+1$      & \\
  \bottomrule
\end{tabular}
\end{center}

\begin{quote}
Wenn alles geklappt hat: \textbf{jedes Element $\ne 0$ hat ein
Inverses}. Das ist genau das, was einen Körper ausmacht.
\end{quote}
\end{bleistiftuebung}

% =============================================================
\section{Python-Werkstatt: GF($2^3$) implementieren \normalfont(\textit{≈ 30 min})}
% =============================================================

Jetzt bauen wir das in Python. Lege eine neue Datei \texttt{gf8.py}
an.

\begin{pythoneinheit}{1}{Die Klasse}
\pythoncodeteil{tag5/gf8.py}{klasse}

\paragraph{Aufgabe 1.1 — Verifiziere mit deiner Bleistift-Tafel.}
Wähle drei Multiplikationen aus Bleistiftübung 5 oder aus deiner
Tafel und prüfe, ob das Python-Ergebnis übereinstimmt.
\end{pythoneinheit}

\begin{pythoneinheit}{2}{Inverse per Brute Force}
Solange wir nur 8 Elemente haben, können wir das Inverse jedes
Elements einfach durch Probieren finden:

\pythoncodeteil{tag5/gf8.py}{inverse}

\paragraph{Aufgabe 2.1 — Vergleiche mit deiner Tafel.}
Stimmen die berechneten Inverse mit Bleistiftübung 7 überein?

\paragraph{Aufgabe 2.2 — Multiplikationstafel automatisch erzeugen.}
Schreibe Code, der die komplette $8 \times 8$-Multiplikationstafel
ausgibt. Format zum Beispiel:

\begin{verbatim}
   ·  | 0  1  2  3  4  5  6  7
   ---+------------------------
   0  | 0  0  0  0  0  0  0  0
   1  | 0  1  2  3  4  5  6  7
   ...
\end{verbatim}

Vergleiche das mit deiner Bleistift-Tafel aus Bleistiftübung 6
(falls du sie gemacht hast).
\end{pythoneinheit}

% =============================================================
\section{Reflexion und Ausblick \normalfont(\textit{≈ 15 min})}
% =============================================================

\subsection*{Was wir heute gelernt haben}

\begin{itemize}
\item \textbf{Körper} sind Zahlensysteme, in denen man uneingeschränkt
  rechnen kann (auch dividieren, außer durch 0).
\item \textbf{GF($p$)} für $p$ prim ist ein Körper — einfach
  Modulo-Rechnung mit Inversen.
\item \textbf{Modulo 256 ist kein Körper} — man kann nicht teilen.
\item Die Lösung: \textbf{Polynom-Arithmetik modulo eines
  irreduziblen Polynoms} liefert Körper für jede Zahl der Form
  $2^n$.
\item \textbf{GF($2^3$)} ist der erste solche Körper, den wir konkret
  aufgebaut haben: 8 Elemente, jedes ein Polynom vom Grad $< 3$,
  Multiplikation modulo $x^3 + x + 1$.
\item \textbf{Jedes Element $\ne 0$ hat ein Inverses} — wir haben die
  Inverse durch Probieren gefunden.
\end{itemize}

\subsection*{Drei Fragen zum Mitnehmen}

\begin{enumerate}
\item Wir haben GF($2^3$) gemacht. Für Bytes brauchen wir GF($2^8$)
  — also 256 Elemente, Polynome vom Grad $< 8$ modulo eines
  irreduziblen Polynoms vom Grad 8. Wie würde sich der Code aus
  Python-Einheit 1 ändern? (Tipp: vier Stellen, alle minimal.)

\item Inverse durch Probieren ist ineffizient. Bei GF($2^8$) wären
  das 255 Versuche pro Inverse. Es gibt einen viel besseren
  Algorithmus (genannt \textbf{Erweiterter Euklidischer
  Algorithmus}). Bevor du davon erfährst — wie würdest \emph{du}
  spontan vorschlagen, das schneller zu machen?

\item Wir haben heute viele Polynome auf einer abstrakten Ebene
  behandelt — als Elemente eines Körpers. Ab morgen werden wir
  Polynome aber wieder als „echte“ Polynome verwenden, deren
  Koeffizienten \emph{Elemente} eines Körpers sind. Klingt
  verwirrend? Es ist genau wie bei reellen Zahlen: $\mathbb{R}$ ist
  ein Körper, und ein Polynom wie $3x^2 + 5x - 2$ hat Koeffizienten
  \emph{aus} $\mathbb{R}$. Wir machen jetzt dasselbe, nur mit
  GF($2^8$) statt $\mathbb{R}$.
\end{enumerate}

\subsection*{Vorschau Tag 6}

Morgen treffen wir auf die zentrale Reed-Solomon-Idee: \textbf{Daten
als Polynom interpretieren, an mehreren Stellen auswerten, und genug
zusätzliche Auswertungen mitspeichern, um aus weniger als der Hälfte
aller Stellen das Original rekonstruieren zu können.} Das ist die
geniale Verallgemeinerung dessen, was wir bei ISBN-10 mit „eine
Ziffer rekonstruieren“ schon gemacht haben.

Dafür brauchen wir die Werkzeuge von heute (Körper-Arithmetik) und
die Werkzeuge von Tag 4 (Polynomdivision). Tag 6 ist also Synthese
der letzten zwei Tage.

